\section{What's new}
\label{sec:novedades}


\begin{itemize}
	\item Next version:
		\begin{itemize}
			\item \tikzname\ patterns in the fills of the stones, the marks, the territory and the background (|fill/pattern|, p.~\pageref{sec:piedras}).
			\item Coordinates on the border in lower case (|label/letter case = lower|, p.~\pageref{subsec:etiquetas}).
			\item The profile |draft|: black and white, with thin lines, no border on the background, and lines where there is a fill.
			\item |goban/arrow| and |goban/line|: the style of the arrows and lines is part of the profile of the board (p.~\pageref{subsec:flechas}). Every profile defines it: |pgfgo| and |pgfgosmooth| draw them in |red!80!black|.
			\item |scale line width|: |scale| also multiplies the width of the lines.
			\item |layout/valign = bottom|: the boards of a row of the grid end at the same height (p.~\pageref{subsec:disposicion}).
			\item |partial| of |\sgfdiagrams|, |\sgfproblem| and |\sgfanimation| also takes two corners (|partial = K10:T19|): the same part in every board (p.~\pageref{subsec:sgfdiagramas}).
			\item |\pgfgodefinesgf| takes |origin| and |transpose|: the corner where the point |aa| of the file goes and whether the first letter is the row, the eight symmetries of the board (p.~\pageref{subsec:sgfnombres}).
			\item |name| in |\sgfdiagrams| and |\sgfproblem|: each diagram is a named board (|var1|, |var2a|...), for |resume| and the views in the text (p.~\pageref{subsec:sgfdiagramas}).
			\item Columns of two letters after |Z|, for boards of more than 25 columns: |AA5|, |partial = AZ5:BG8| (p.~\pageref{subsec:sintaxis-de-coordenadas}).
			\item |links| in |\sgfproblem|: with |hyperref|, the statement and its solutions link to each other (p.~\pageref{subsec:sgfproblemas}).
			\item The option |pages|: one board per page, and |<document>-pages.lua| turns them into SVG, PNG or PDF (p.~\pageref{subsec:paginas}).
			\item Sensei's Library and L19 diagrams (|$$| lines), in a file or in the environment |pgfgodiagram|: every SGF command reads them (p.~\pageref{subsec:sgfsensei}).
			\item |goban/blank|: intersections that are not part of the board, without lines or star points; they are the |_| of Sensei's Library diagrams (p.~\pageref{subsec:parcial}).
			\item The sizes of the board are in |pt|: they do not change with the size of the font of the document. One character of |mark/sequence| and one of |mark/custom| have the same size.
			\item |grid/edge|: the edge of the grid, thicker or in another color.
			\item |layout/align| is renamed |layout/halign|, and |layout/align captions| puts the captions of a row on the same line with the boards at the top (p.~\pageref{subsec:disposicion}).
			\item |notes/font|: the font of the notes.
			\item |\sgfproblem| leaves out the games that the file does not have, with a single warning.
			\item In Spanish, the colors are singular and the result also agrees with a name: ``Gana Negro'', ``Juega Blanco'', ``Gana Ana por abandono''. In Portuguese, ``Vitória de Ana''.
		\end{itemize}
	\item \pgfgoversion:
		\begin{itemize}
			\item Kernel and rules of the game:
				\begin{itemize}
					\item The computations (coordinates, numbering of moves, memory, etc.) have been separated from the drawing. They are done by a \emph{kernel}, chosen with the package option |kernel| (p.~\pageref{subsec:kernel}). Besides the kernel written in L3, a kernel written in Lua has been added, which is the default with Lua\LaTeX.
					\item With the Lua kernel, captures have been added (|captures|, p.~\pageref{subsec:capturas}). With |captured marks = remove|, the marks on a captured stone go with it, and |\pgfgoprisoners| counts the prisoners of each player.
					\item ``N at M'' notes have been added for numbered moves on points already played (|notes| and |\pgfgomovenotes|, p.~\pageref{subsec:notas}). The numbered passes of an SGF are noted as ``N: pass'' (|notes/pass|).
					\item |\territory|, which marks the territory (with a polygon, or by itself with |auto| and the Lua kernel), and |\pgfgoscore|, which counts the points (p.~\pageref{subsec:territorio}), have been added. With |draw = region|, the territory of each color is a single polygon, and |\pgfgoregion| gives its measures.
				\end{itemize}
			\item Named boards and views:
				\begin{itemize}
					\item A |goban| with |name| keeps everything it draws, and another one starts with that board with |resume| (p.~\pageref{subsec:recordar}). |resume| brings the stones; |numbers|, |marks| and |territory| bring the rest. With |name| and |resume| in the same |goban|, variations are built.
					\item |\pgfgoview| draws a point or part of a named board in the middle of the text (p.~\pageref{subsec:vistas}), with the options of |view| (|marks|, |arrows|, |lines|, |numbers|, |territory|, |forgotten|, |raise|, |profile|).
					\item A view can come before its |goban|, like |\ref| (p.~\pageref{subsec:vistas-antes}), and a board can be used in another document with |export| and |\pgfgoimport| (p.~\pageref{subsec:exportar}).
				\end{itemize}
			\item SGF:
				\begin{itemize}
					\item With the Lua kernel, |\sgf| reads games in SGF (p.~\pageref{subsec:sgftablero}): variations (|variation|), marks of the file (|marks|), the numbering of |MN| and |restart|, and the size of the board taken from the file (|sgf size|).
					\item Collections of games (|game|), the view (|VW|, |sgf view|), figures (|FG|) in |\sgfdiagrams|, variations as letters (|variations|), |PL| and files in Latin-1.
					\item The numbering asked for by |PM| (|numbers| and |restart = auto|), and the view (|VW|) when it changes during the game (|sgf view move|; |\sgfdiagrams| and |\sgfanimation| use it by themselves). |step| accepts a list, and |restart = {from = |\meta{n}|}| numbers each stretch from \meta{n}.
					\item |variation = {R16, P17, pass}|: the variation can also be given with its moves, from the first one.
					\item Files in GBK, Big5, Shift\_JIS, EUC-JP and EUC-KR are read (|CA| or |sgf charset|). If the font lacks the Chinese, Japanese or Korean characters of a text of the file, the package warns.
					\item Names for the files, with |\pgfgodefinesgf| (p.~\pageref{subsec:sgfnombres}): the path, the game and the charset are written only once.
					\item The information of the game (|\sgfinfo|, also of a move with |move|). The names of the nodes and the annotations (|TE|, |BM|, |GB|, etc.) are written with the comments. |\sgfcomment| writes the comment of a move, and |chat = false| removes the chat of the servers (p.~\pageref{subsec:sgfdatos}).
					\item |\sgfdiagrams| draws a game as a series of diagrams, as in books (p.~\pageref{subsec:sgfdiagramas}). It writes the comments below each diagram, beside it or after all of them (|comments|), and the result, in words, and the score below the last one (|result|). With |variations|, it also draws the variations, each in its own diagram.
					\item |last mark| marks the last move, and |to = last| goes to the end.
					\item |\sgfproblem| has been added, which draws a problem (tsumego): the statement and one diagram for each variation of the solution, with its comments (p.~\pageref{subsec:sgfproblemas}). With |game = 1:20|, a problem sheet from a collection; with |verdict|, the caption says whether the variation is right.
					\item With the |animate| package, |\sgfanimation| steps through a game with buttons (p.~\pageref{subsec:sgfanimacion}). It builds the animation in layers (|layers|): each stone is drawn only once, and the PDF of a long game takes much less space.
				\end{itemize}
			\item Drawing:
				\begin{itemize}
					\item Arrows and lines have been added (|\arrow| and |\line|, p.~\pageref{subsec:flechas}), with |dash|, |shorten|, |bend| and |tips|, and the mark |dim|. |\sgf| also draws |AR|, |LN|, |SL| and |DD|.
					\item The board draws the star points (|hoshi|, p.~\pageref{subsec:hoshi}).
					\item The coordinates of the border can be letters, numbers or kanji on each axis, counted from either side (|columns|, |rows|, |style = japanese|, etc.).
					\item |partial = auto| shows only the part of the board used by what is drawn (p.~\pageref{subsec:parcial}); also in |\sgfdiagrams|, |\sgfanimation| and |\sgfproblem|.
					\item Images (PNG, JPG or PDF) on the stones (|stone/image|, p.~\pageref{subsec:piedrasimagen}), on the background and the grid of the board, on the star points and on the marks. With several files, each point always uses the same one, and |rotate| turns each stone or tile. Each image is kept only once in the PDF (p.~\pageref{subsubsec:imagenespdf}).
					\item Shadows and shines of your own (|shadow/image| and |shine/image|), and |shadow/scale| and |shadow/color|.
					\item Marks of your own, with PGF code or with an image (|\pgfgodefinemark|, p.~\pageref{subsec:marcaspropias}), and |mark/scale|.
					\item Numbers in kanji for the moves and the marks (|format = kanji|).
					\item Opacity of the stroke and of the fill of the marks, of the numbers of the moves, of the images of the stones and of the star points.
					\item The \tikzname\ coordinate system |go|: |(go cs:C3)| (p.~\pageref{subsec:tikz}).
				\end{itemize}
			\item Styles and profiles:
				\begin{itemize}
					\item Profiles are \emph{keys} (p.~\pageref{subsec:perfiles}): |profile|, |goban profile|, |players profile| and |players/|\meta{player}|/profile|. Before they are applied, the board and the players go back to the initial style. |\pgfgodefineprofile| defines a whole profile at once.
					\item Three new profiles: |simple|, and |shell| and |yunzi|, with images (p.~\pageref{subsec:perfilshell}): a wooden board with slate and shell stones, or matte ones. The images come with the package, next to its |.sty| files, and are used by their name.
					\item The same rule for the \emph{keys} that turn something on or off (p.~\pageref{subsec:prender}).
					\item The options of a |goban| accept those of |\pgfgoset|, only for that board, and |\pgfgoset{mark = {...}}| also sets the style of the marks of the three players.
				\end{itemize}
			\item Layout: |layout| and the |pgfgolayout| environment lay out the diagrams in a grid of columns, with |fit|, |width = natural|, |per page| and |per frame| (p.~\pageref{subsec:disposicion}).
			\item Languages: the texts of the package are translated into ten languages and follow |babel| or |polyglossia|; |lang|, |\pgfgotext|, |\pgfgodefinelanguage| and |\pgfgolanguagealias| (p.~\pageref{subsec:idiomas}).
			\item Coordinates:
				\begin{itemize}
					\item Lines in any direction and stairs: |A1 -- E6| (in any direction), \verb+A1 -|- E6+ and \verb+A1 |-| E6+.
					\item Rectangles are written |A1:C3|, as in |partial|, and stretches and rectangles are chained with the operators.
					\item |\moves| accepts every form of |\stones| and plays each one in its order (p.~\pageref{subsec:orden-coordenadas}); |\moves| and |\marks[sequence]| warn if they have no fixed order.
				\end{itemize}
			\item Behavior changes (the full list, with the old names of the \emph{keys}, is in |CHANGELOG.md|):
				\begin{itemize}
					\item |\moves| numbers its moves; |numbers = false| leaves them without a number. |\move|, |\move*| and |\moves*| have been removed.
					\item |\stones| without a player draws neutral stones, like |\stone|.
					\item The only way to remember a board is |name| and |resume|; |goban/remember|, |goban/forget| and |resume| without a name have been removed.
					\item |\pgfgouseprofile| and the other |\pgfgouse...| have been removed; now |\pgfgoset{profile = wain}| is used.
					\item The operator |R| has been removed (|A1 R C3| is |A1:C3|), and |:--:|, \verb+:-|:+ and \verb+:|-:+ are now |--|, \verb+-|-+ and \verb+|-|+.
					\item The territory is drawn when the |goban| ends, with the last |\territory| that marked each point.
					\item A move number is never covered by a label of the file (|LB|) or by |last mark|.
				\end{itemize}
		\end{itemize}
	\item 0.2025.10.13: Bug fixes. Each |goban| starts at move 1 with Black, |\mark| and |\marks| only exist inside |goban|, and what falls outside a partial board is no longer drawn. The full list is in the file |CHANGELOG.md|.
	\item 0.2025.03.10:
		\begin{itemize}
			\item The \emph{keys} |inner| and |outer| have been added for the \emph{rounded corner} of the board.
			\item Three new coordinate operators have been added: |*| for every intersection, |L| for the whole column of coordinate L and |N| for the whole row of coordinate N.
		\end{itemize}
	\item 0.2025.03.07: First release.
\end{itemize}
